\documentclass{article}
\usepackage[UTF8]{ctex}
\usepackage{amsmath,amssymb,amsfonts}
\usepackage{enumitem}
\usepackage{geometry}
\geometry{a4paper, margin=1in}
% === 错题本排版优化（不参与题目标识） ===
\setlength{\parindent}{0pt}
\setlength{\parskip}{0.35em plus 0.10em minus 0.08em}
\setlength{\emergencystretch}{3em}
\setlist{itemsep=0.25em, topsep=0.40em, parsep=0pt}
\allowdisplaybreaks[2]
% === 错题本排版优化结束 ===
\begin{document}
\section*{山东大学2024年高等代数真题}
\section*{113．山东大学2024年线性代数与常微分方程试卷}
一．高等代数。\\

\subsection{题目1}
1．（15 分）考虑齐次线性方程组

\[
\left\{\begin{array}{c}
a x_{1}+b x_{2}+b x_{3}+\cdots+b x_{n}=0 ; \\
b x_{1}+a x_{2}+b x_{3}+\cdots+b x_{n}=0 ; \\
\cdots \cdots \\
b x_{1}+b x_{2}+b x_{3}+\cdots+a x_{n}=0 .
\end{array}\right.
\]

其中 \(a \neq 0, b \neq 0, n \geq 2\) 。试讨论 \(a, b\) 取何值时，方程组仅有零解？有无穷多解？并在有无穷多解时，用基础解系给出其通解。\\

\subsection{题目2}
2．（15 分）设 \(A, B, C\) 分别为 \(m \times n, n \times t, s \times m\) 阶矩阵。\\
（1）若矩阵 \(A\) 的秩 \(r(A)=r\) ，证明：存在可逆阵 \(P, Q\) ，使得 \(P A\) 的后 \(m-r\) 行全为零，\(A Q\) 的后 \(n-r\)列全为零。\\
（2）利用（1）证明：若 \(r(A)=n\) ，则 \(r(A B)=r(B)\) ；若 \(r(A)=m\) ，则 \(r(C A)=r(C)\) 。\\

\subsection{题目3}
3．（10 分）设 \(A\) 为 \(n\) 阶正定矩阵，\(\alpha\) 为 \(n\) 维实的列向量，证明：\(A^{-1}\) 与 \(A+\alpha \alpha^{\mathrm{T}}\) 均为正定矩阵，其中\(A^{-1}\) 为 \(A\) 的逆矩阵，\(\alpha^{\mathrm{T}}\) 为 \(\alpha\) 的转置。\\

\subsection{题目4}
4．（15 分）设 \(A=\left(\begin{array}{ccc}0 & b & -c \\ -b & 0 & a \\ c & -a & 0\end{array}\right)\) 为实矩阵，令 \(B=A^{2}+q A+E\) ，其中 \(q=a^{2}+b^{2}+c^{2}, E\) 为三阶单位阵。试问：当且仅当 \(q\) 为何值时，矩阵 \(B\) 是正交矩阵？\\

\subsection{题目5}
5．（15 分）设 \(A, B\) 为3阶复方阵，且都只有一个特征值 \(\lambda_{0}\) 。证明：\(A\) 与 \(B\) 相似的充要条件是

\[
\operatorname{dim}\left(V_{\lambda_{0}}(A)\right)=\operatorname{dim}\left(V_{\lambda_{0}}(B)\right) .
\]

\subsection{题目6}
6．（20 分）设 \(V\) 是 \(n\) 维欧氏空间，\(W_{1}, W_{2}\) 是 \(V\) 的子空间，且 \(\operatorname{dim} W_{1}=s<\operatorname{dim} W_{2}=t\) 。证明：\\
（1）存在 \(\beta \in W_{2}, \beta \neq 0\) ，而 \(\left(\beta, W_{1}\right)=0\) ，且 \(\operatorname{dim}\left(W_{1}^{\perp} \cap W_{2}\right) \geq t-s\) 。\\
（2） \(\operatorname{dim}_{\text {d微 }}\left(W_{\text {分 }}\right.\) 程 \(\left.W_{2}^{\perp}\right) \leq n-t+s\) ．\\

\subsection{题目7}
1．（10 分）求方程 \(\left(y+x^{3} y+2 x^{2}\right) \mathrm{d} x+\left(x+4 x y^{4}+8 y^{3}\right) \mathrm{d} y=0\) 的通解。\\

\subsection{题目8}
2．（10 分）求方程 \(\left(y^{\prime}\right)^{3}+y^{3}-3 y y^{\prime}=0\) 的通解．\\

\subsection{题目9}
3．（10 分）试证若 \(y=\varphi(x)\) 是方程 \(\frac{\mathrm{d} y}{\mathrm{~d} x}=p(x) \sin y\) 的满足初试条件 \(\varphi(0)=0\) 的解，则 \(\varphi(x) \equiv 0\) ，其中 \(p(x)\) 在 \((-\infty,+\infty)\) 上连续。\\

\subsection{题目10}
4．（20 分）设 \(t>0, x, y\) 是关于 \(t\) 的函数，解方程组

\[
\left\{\begin{array}{l}
t x^{\prime}-x-y=0 \\
t y^{\prime}+x-y=0
\end{array}\right.
\]

\subsection{题目11}
5．（10 分）是否存在 \(\mathbb{R}\) 上连续函数 \(p, q\) ，使得微分方程 \(y^{\prime \prime}+p(x) y^{\prime}+q(x) y=0, x \in \mathbb{R}\) 有两个解\(\phi(x)=\sin x, \psi(x)=x e^{x}, x \in \mathbb{R}\) ？
\end{document}
